% =============================================================================
% Lecture 2: System Paradigms
% Introduction to Machine Learning Systems (Volume I)
% =============================================================================
% Physics sets four deployment regimes. The lecture teaches the power
% ladder Lecture 0 showed as a teaser, then the four physical constraints
% behind it (distance, memory bandwidth, power and heat, moving data), with
% the chapter's napkin math as worked examples. The roofline appears here
% first. It closes on choosing a tier and splitting work across tiers.
%
% Build: make -C slides build-02_system_paradigms   (figures: images/make_figs.py)
% Plan:  lecture_map.yml in this folder.
% =============================================================================

\documentclass[aspectratio=169, 12pt]{beamer}
\usepackage{../../assets/beamerthememlsys}
\mlsyssetup{
  volume={Volume I},
  chapter={Chapter 2},
  chaptertitle={System Paradigms},
  logo={../../assets/img/logo-mlsysbook.png},
  instlogo={../../assets/img/logo-harvard.png},
}
\usepackage[T1]{fontenc}
\usepackage[scaled=0.9]{helvet}
\usepackage{courier}
\renewcommand{\familydefault}{\sfdefault}
\usepackage{booktabs,amsmath}
\usetikzlibrary{arrows.meta,positioning}
% a stage box in the hybrid-pattern figure, colored by tier (RedLine/RedL, ...)
\tikzset{bx/.style={draw=#1Line, fill=#1L, line width=0.8pt, minimum width=2.3cm, minimum height=0.52cm, inner sep=1pt, font=\footnotesize, text=darktext}}
% show a TikZ scope only on some overlay steps, keeping its space on the others
\tikzset{invisible/.style={opacity=0, text opacity=0}, visible on/.style={alt={#1{}{invisible}}}, alt/.code args={<#1>#2#3}{\alt<#1>{\pgfkeysalso{#2}}{\pgfkeysalso{#3}}}}
\graphicspath{{images/}}
\setcounter{mlsystotalsections}{5}

\title{System Paradigms}
\author{Vijay Janapa Reddi}
\institute{Harvard University}
\date{}

% A large answer choice for predict and concept-question slides
\newcommand{\choice}[2]{\textbf{\textcolor{mlsysaccent}{#1}}\enspace #2}
% A highlighted (revealed) answer choice
\newcommand{\rchoice}[2]{\colorbox{mlsysaccent!15}{\choice{#1}{#2}}}
% A photo with its tier label underneath
% A takeaway thumbnail of the figure that taught the line
\newcommand{\thumb}[1]{\safeimg[width=3.1cm,height=0.95cm,keepaspectratio]{#1}}
\newcommand{\tierphoto}[3]{\centering\parbox[c][3.2cm][c]{\linewidth}{\centering#1}\\[0.08cm]{\footnotesize\textbf{#2}}\\{\scriptsize\textcolor{midgray}{#3}}}

\begin{document}

% =============================================================================
% TITLE
% =============================================================================
\mlsysnote{[1 min] Read the thesis aloud once. Lecture 0 showed a ladder of four machines, from a 3 MW datacenter to a 1 mW microcontroller, and promised to explain it. Today does that: each rung exists because a physical law binds there. The notes total about 71.5 minutes.}
\mlsystitle{System Paradigms}{Physics sets four deployment regimes: distance, memory bandwidth, power, and the cost of moving data decide where a model can run.}{cover_system_paradigms_blueprint.png}

% =============================================================================
% HOOK
% =============================================================================
\begin{frame}{Predict: will the arm stop in time?}
  \mlsysactivity{Predict}{3 min}
  \begin{columns}[c]
    \begin{column}{0.50\textwidth}
      \centering\safeimg[width=\textwidth,height=4.9cm,keepaspectratio]{robot-arms-bmw-leipzig.jpg}
    \end{column}
    \begin{column}{0.47\textwidth}
      \normalsize A camera must stop the arm within \textbf{10 ms} of seeing a person. The model runs in a cloud region \textbf{1,500 km} away.\\[0.4cm]
      \large
      \choice{A}{Yes}\\[0.12cm]
      \choice{B}{No}\\[0.12cm]
      \choice{C}{Depends on the cloud GPU}
    \end{column}
  \end{columns}
  \mlsyscredit{Photo: BMW Werk Leipzig (CC BY-SA 2.0 de), Wikimedia Commons. Scenario: System Paradigms, Napkin Math 2.2.}
  \note{[3 min] The photo is a car-body line; workers and robots share the floor. A safety monitor watches a camera and must stop an arm within 10 ms of seeing a person. The model runs in a regional cloud datacenter 1,500 km away. ENGAGE: vote with fingers, one to three, and ask two students to defend their vote. Do not reveal the answer; the Distance segment returns to it with the arithmetic. Expect many votes for C: students picture the GPU as the slow part. That belief is what today corrects. Also write down the share who chose A, to show the room at the reveal.}
\end{frame}

\begin{frame}{Recall: what are the iron law's three terms?}
  \mlsysactivity{Retrieval}{1.5 min}
  \centering
  {\Large\bfseries Close your notes.}\\[0.3cm]
  {\normalsize Write the time for one inference as a sum of three terms.}\\[0.6cm]
  \uncover<2->{\Large$T = \underbrace{\dfrac{\Dvol}{\text{BW}}}_{\text{data movement}} + \underbrace{\dfrac{O}{\Rpeak \cdot \eta_{\text{hw}}}}_{\text{compute}} + \underbrace{\Llat}_{\text{overhead}}$}
  \mlsyscredit{Lecture 1; System Paradigms, equation 4.}
  \note{[1.5 min] Forty-five seconds of silent writing, then step 2. Data movement is the bytes moved divided by bandwidth; compute is the operations divided by the rate the hardware actually achieves; overhead is fixed latency, which today includes the network. LINK: Lecture 1 built this equation. Today each term gets a physical law that sets its size on each kind of machine.}
\end{frame}

\begin{frame}{Why does the same model need different engineering on a phone?}
  \normalsize By the end of today, you will be able to:\\[0.25cm]
  \begin{itemize}\setlength{\itemsep}{0.25cm}
    \item Explain how distance, memory bandwidth, power, and data volume create four deployment regimes.
    \item Place a workload on a roofline and name the resource that limits it.
    \item Estimate latency, bandwidth, and energy budgets with napkin math.
    \item Choose a tier, or a split across tiers, for an application.
  \end{itemize}
  \note{[1.5 min] The title is the driving question; the full form is ``why does the same model need different engineering on a phone and in a datacenter?'' Read the four outcomes. The third is the habit this course drills: a few lines of arithmetic before any code. FLEX: [CORE].}
\end{frame}

% =============================================================================
% SECTION: FOUR REGIMES
% =============================================================================
\mlsysnote{[0.5 min] Start with the four machines themselves, then what each one offers a model.}
\mlsysdivider{Why are there four kinds of ML machine, not one?}{System Paradigms, \S2.1 and \S2.4}
\section{Four Regimes}

\begin{frame}{The same task can run on four very different machines}
  \begin{columns}[T]
    \begin{column}{0.24\textwidth}
      \tierphoto{\safeimg[width=\textwidth,height=3.1cm,keepaspectratio]{book-cloud-datacenter.jpg}}{Cloud}{a datacenter}
    \end{column}
    \begin{column}{0.24\textwidth}
      \uncover<2->{\tierphoto{\safeimg[width=\textwidth,height=3.1cm,keepaspectratio]{jetson-agx-orin.jpg}}{Edge}{a box beside the sensors}}
    \end{column}
    \begin{column}{0.24\textwidth}
      \uncover<3->{\tierphoto{\mlsysphoto[width=\textwidth,height=3.1cm,keepaspectratio]{smartphone-in-hand.jpg}}{Mobile}{a phone}}
    \end{column}
    \begin{column}{0.24\textwidth}
      \uncover<4->{\tierphoto{\safeimg[width=\textwidth,height=3.1cm,keepaspectratio]{esp32-s3-board.jpg}}{TinyML}{a microcontroller}}
    \end{column}
  \end{columns}
  \vspace{0.35cm}
  \centering\uncover<4->{\normalsize Classifying one image is the same task on all four.}
  \mlsyscredit{Photos: Google, via System Paradigms fig.\ 5; Auledas (CC BY 4.0); Dennis Cort\'es (CC0); VectorVoyager (CC BY-SA 4.0); Wikimedia Commons.}
  \note{[2 min] Real things first. Step 1: a Google datacenter aisle, rows of accelerator racks. Step 2: an NVIDIA Jetson AGX Orin, an edge computer that sits in a store, a factory, or a car, next to the cameras it serves. Step 3: a phone. Step 4: an ESP32-S3 board, a microcontroller that costs under ten dollars. The chapter calls these four regimes cloud, edge, mobile, and TinyML. WARN: edge is a location, not a hardware class; it means ``near the data,'' and spans gateways to on-premises servers. ENGAGE: ask which of the four is in the room right now (phones, and probably microcontrollers inside every phone and laptop). FLEX: [CORE].}
\end{frame}

\begin{frame}{Predict: how much power does each tier draw?}
  \mlsysactivity{Predict}{3 min}
  \centering
  \mlsyssteps[width=0.95\textwidth,height=5.0cm,keepaspectratio]{power_tiers_step}{4}\\[0.1cm]
  \uncover<4->{\normalsize Power spans nine orders of magnitude; memory falls 250-million-fold.}
  \mlsyscredit{System Paradigms, \S2.1 and table 4. Memory is that of the example platform: TPU v4 pod, DGX Spark, flagship phone, ESP32-S3 SRAM.}
  \note{[3 min] ENGAGE before step 1: students write a guess in watts for each of the four machines (30 seconds). Most guess the phone too high and the microcontroller far too high; few guess a datacenter pod in megawatts. Explain the axis first: power on a log scale, each grid line a factor of 100. Step 1: a cloud pod of 4,096 TPU v4 chips draws about 3 MW and holds 128 TB of high-bandwidth memory. Step 2: an edge computer, about 200 W, with 128 GB (the chapter's example is an NVIDIA DGX Spark, larger than the Jetson in the photo). Step 3: a phone, about 5 W, with 8 to 16 GB. Step 4: a microcontroller averages under 1 mW when duty-cycled (awake a few milliseconds, asleep otherwise), with 512 KB of SRAM, the small, fast memory on the chip itself. Megawatts to milliwatts is nine orders of magnitude. ENGAGE: ``Where would your laptop sit?'' (Between edge and mobile: tens of watts, more than a phone, less than an edge server.) WARN: the cloud number is a whole pod, the others single devices; the chapter compares deployment scales, not equivalent chips. FLEX: [CORE].}
\end{frame}

\begin{frame}{Each tier offers a model a different silicon contract}
  \centering\footnotesize
  \renewcommand{\arraystretch}{1.35}
  \begin{tabular}{@{}P{0.27\textwidth}P{0.22\textwidth}P{0.22\textwidth}P{0.20\textwidth}@{}}
    \toprule
    \textbf{Tier and example} & \textbf{Tier compute} & \textbf{Tier memory bandwidth} & \textbf{Memory} \\
    \midrule
    Cloud: TPU v4 pod & over 1,000 TFLOP/s & over 1,000 GB/s & 128 TB (pod) \\
    Edge: DGX Spark & about 1 PFLOP/s (sparse FP4) & over 270 GB/s & 128 GB \\
    Mobile: flagship phone & 15 to 45 TOPS & 51.2 to 77 GB/s & 8 to 16 GB \\
    TinyML: ESP32-S3 & under 1 TOPS & --- & \colorbox{mlsysaccent!15}{512 KB SRAM} \\
    \bottomrule
  \end{tabular}\\[0.1cm]
  {\scriptsize\textcolor{midgray}{Compute and bandwidth are per device for every tier; the whole 4,096-chip pod delivers 1.1 EFLOP/s.}}\\[0.3cm]
  \small Compute, bandwidth, and overhead set the time; memory and energy decide whether it runs.
  \mlsyscredit{System Paradigms, table 4 (per-device tier ranges). Point-in-time products; the reasoning outlasts them.}
  \note{[2 min] The chapter calls these terms the silicon contract: every platform offers a compute rate, a memory bandwidth, and a fixed overhead, and those three enter the iron law. Units first: a TFLOP/s is $10^{12}$ floating-point operations per second; memory bandwidth is the bytes per second memory delivers to the arithmetic units. The table shows memory in place of overhead, because overhead (network round trips, launch costs) depends on the path, not the chip; the Distance segment measures it. Compute and bandwidth are per-device ranges, as in the chapter; only the cloud memory is the whole pod's. Walk one row: a phone NPU (neural processing unit, the phone's ML accelerator) offers 15 to 45 trillion operations per second (TOPS, integer operations) but only 51 to 77 GB/s from its memory. The edge figure is quoted in sparse FP4: 4-bit numbers, counted as if half the weights were zero, so it overstates the dense rate a typical model gets. Two more terms are not in the time equation at all: memory capacity and energy. They are preconditions. A model whose working set does not fit in 512 KB of SRAM never runs, however long you wait. WARN: students read TOPS as the whole story; bandwidth often binds, as the Memory segment shows. FLEX: [CORE].}
\end{frame}

\mlsysnote{[1.5 min] The key law for today. The iron law from Lecture 1, with the physical wall that sets each term on each tier. Data movement: the memory wall, compute outgrowing bandwidth. Compute: the power wall, because heat caps how fast a chip can run. Overhead: the light barrier, because a remote request pays the round trip. Below the equation: capacity and energy are checked first; a model that does not fit, or a battery that cannot pay, ends the discussion before any term is evaluated. The next three segments take one wall each. For students who know computer architecture: these are the classic walls; what ML changes is the scale of data each request moves, which the next segment shows.}
\mlsysfocus{Three walls set the iron law}{%
  $\displaystyle T = \underbrace{\frac{\Dvol}{\text{BW}}}_{\text{memory wall}} +
   \underbrace{\frac{O}{\Rpeak \cdot \eta_{\text{hw}}}}_{\text{power wall}} +
   \underbrace{\Llat}_{\text{light barrier}}$\\[0.6cm]
  {\small\textcolor{white!75}{Checked first: does it fit in memory, and can the energy budget pay for it?}}%
}

% =============================================================================
% SECTION: DISTANCE
% =============================================================================
\mlsysnote{[0.5 min] The first wall: distance. Return to the arm.}
\mlsysdivider{What does it cost to move a request, or its data?}{System Paradigms, \S2.2, \S2.5, and \S2.6}
\section{Distance}

\begin{frame}{Light in fiber covers only 200 km per millisecond}
  \centering
  {\normalsize$\displaystyle \text{Round trip}_{\min} = \frac{2 \times \text{distance}}{200{,}000\ \text{km/s}}$}\\[0.25cm]
  \mlsyssteps[width=0.92\textwidth,height=3.9cm,keepaspectratio]{rtt_distance_step}{2}
  \mlsyscredit{System Paradigms, equation 1 and \S2.2. Real networks add routing, queuing, and processing on top.}
  \note{[2 min] Light in glass travels at about 200,000 km/s, two thirds of its speed in vacuum. A request goes out and the answer comes back, so the distance counts twice. Step 1: a regional datacenter 1,500 km away costs 15 ms of pure propagation. The green line is the arm's 10 ms budget. Step 2: California to Virginia, about 3,600 km, costs about 36 ms. These are floors: real networks add routing, protocol handshakes, and queuing, so measured round trips are larger. WARN: students think 5G or faster links fix this. Bandwidth improves; propagation does not. No engineering makes light faster. FLEX: [CORE].}
\end{frame}

\begin{frame}{The arm's cloud path misses its budget before any compute}
  \begin{columns}[T]
    \begin{column}{0.40\textwidth}
      \vspace{0pt}\begin{mlsyscard}{mlsysaccent}
        \textbf{Given}\\[0.1cm]
        Safety budget: 10 ms\\
        Distance to the cloud: 1,500 km\\
        Light in fiber: 200,000 km/s
      \end{mlsyscard}
    \end{column}
    \begin{column}{0.52\textwidth}
      \vspace{0pt}\normalsize
      \uncover<2->{Round trip $= \dfrac{2 \times 1{,}500\ \text{km}}{200{,}000\ \text{km/s}} = 15\ \text{ms}$\\[0.35cm]}
      \uncover<3->{15 ms $>$ 10 ms, with zero time left for the model.\\[0.3cm]}
      \uncover<4->{{\Large\textbf{\textcolor{mlsysaccent}{Answer: B. No GPU can fix it.}}}\\[0.3cm]}
      \uncover<5->{Move the model beside the arm; then check the local path.}
    \end{column}
  \end{columns}
  \mlsyscredit{System Paradigms, Napkin Math 2.2 (the distance penalty).}
  \note{[2 min] Back to the hook. Step 2: plug in. Step 3: the round trip alone exceeds the budget by 5 ms before the camera frame is even processed. Step 4: the answer is B. Show the room its vote. Students who chose C pictured a slow GPU; even an infinitely fast cloud GPU leaves 15 ms. This is what the chapter means by a regime: optimizing inside an infeasible regime cannot help. Step 5: running the model on an edge computer beside the arm removes the propagation term, but the local sensing, inference, and stop command must still be measured against 10 ms. LINK: the Choosing segment calls this the latency gate. FLEX: [CORE].}
\end{frame}

\begin{frame}{100 raw camera feeds need 15 times the factory's link}
  \centering
  \mlsyssteps[width=0.95\textwidth,height=4.2cm,keepaspectratio]{camera_bandwidth_step}{3}\\[0.15cm]
  \uncover<3->{\normalsize Defect records need 20 KB/s; raw video needs 18.7 GB/s.}
  \mlsyscredit{System Paradigms, Napkin Math 2.4: 100 cameras, 1080p, 30 FPS, 3 bytes per pixel.}
  \note{[2 min] Distance is not the only cost of moving data; volume is the other. A factory runs quality control on 100 cameras at 1080p, 30 frames per second. Step 1: one camera produces 1920 times 1080 pixels times 3 bytes times 30 frames, about 186.6 MB/s; 100 cameras produce about 18.7 GB/s. Step 2: a dedicated 10 Gb/s fiber line carries 1.25 GB/s (divide bits by 8). The raw stream needs about 15 times more. Step 3: an edge server that runs the detector locally sends only defect records, about 1 KB per detection at about 20 per second: 20 KB/s, a reduction of about 930,000 times. WARN: the most common error is comparing GB/s with Gb/s; ask a student to say where the factor of 8 went. Compression helps but must keep the detail the defect detector needs. FLEX: [CORE].}
\end{frame}

\begin{frame}{Estimate: can a drone offload obstacle avoidance?}
  \mlsysactivity{Your turn}{3 min}
  \begin{columns}[T]
    \begin{column}{0.44\textwidth}
      \vspace{0pt}\begin{mlsyscard}{mlsysaccent}
        \textbf{Given}\\[0.1cm]
        One 4K frame: 24.9 MB\\
        Uplink: 100 Mb/s\\
        Remote response: 110 ms\\
        Deadline: 16.7 ms (60 FPS)
      \end{mlsyscard}
    \end{column}
    \begin{column}{0.52\textwidth}
      \vspace{0pt}\normalsize
      \uncover<2->{Send: $\dfrac{24.9\ \text{MB} \times 8}{100\ \text{Mb/s}} = 1{,}990.7\ \text{ms}$\\[0.25cm]}
      \uncover<3->{Remote path $\geq$ 1,990.7 + 110 = 2,100.7 ms\\[0.25cm]}
      \uncover<4->{{\large\textbf{\textcolor{mlsysaccent}{About 126 times the deadline}}}}
    \end{column}
  \end{columns}
  \vspace{0.1cm}
  \uncover<5->{\centering\small Data locality invariant:\enspace
    $T_{\text{remote}} \ge \underbrace{\textcolor{RedLine}{\Dvol / \text{BW}_{\text{net}}}}_{\text{send: no GPU helps}} + \underbrace{L_{\text{net}}}_{\text{round trip}} + \underbrace{O / (R_{\text{peak,remote}}\,\eta_{\text{hw}})}_{\text{remote compute}}$\\[0.05cm]
    If this lower bound exceeds the deadline, reject the remote path.\par}
  \mlsyscredit{System Paradigms, Napkin Math 2.6 (the locality crossover). Uncompressed frames, as in the chapter.}
  \note{[3.5 min] ENGAGE: one minute with a neighbor before step 2. Step 2: convert bytes to bits, then divide. Sending one frame takes about two seconds. Step 3: add the remote response; this is a lower bound because it ignores queuing and serialization. Step 4: about 126 times the 16.7 ms frame deadline. A faster cloud GPU shrinks only the 110 ms. Step 5: the chapter's data locality invariant (Definition 2.3). Read the three terms left to right: the time to send the input over the network, the network round trip, and the remote compute. The first term, in red, depends only on the bytes and the link; a faster GPU cannot touch it. If even this lower bound exceeds the deadline, reject the remote path; compare measured times only among the paths that survive. Here the send term alone is 1,990.7 ms, so no remote design can pass. The same arithmetic at dataset scale is data gravity: moving a large dataset to a cluster eventually costs more than moving the code to the data. WARN: a student may object that video is compressed; it is, but the conclusion survives a 100-fold compression (still about 20 ms of sending alone). FLEX: [CORE].}
\end{frame}

% =============================================================================
% SECTION: MEMORY
% =============================================================================
\mlsysnote{[0.5 min] The second wall lives inside every machine: moving bytes from memory to the arithmetic units.}
\mlsysdivider{Why do fast processors sit idle waiting for bytes?}{System Paradigms, \S2.2, \S2.3, and \S2.4}
\section{Memory}

\begin{frame}{Compute has outgrown memory bandwidth for decades}
  \centering
  \safeimg[width=0.92\textwidth,height=5.0cm,keepaspectratio]{memory_wall.pdf}
  \mlsyscredit{System Paradigms, \S2.2, equation 3 (illustrative rates); Wulf and McKee 1995.}
  \note{[2 min] Axes: years on the bottom, growth on a log scale. Orange: peak compute doubling about every 18 months, about 1.6 times a year. Blue: memory bandwidth growing about 20 percent a year. After a decade compute is about 102 times faster and bandwidth about 6 times faster, so compute is about 16 times further ahead. Wulf and McKee named this the memory wall in 1995. The consequence: each hardware generation makes data movement a larger share of time and energy, unless the workload reuses each byte more. WARN: these are the chapter's illustrative rates, not a forecast for any one product line. For students who know architecture: the same wall that motivated caches; ML adds workloads that touch every weight for every token. FLEX: [OPTIONAL]; the roofline carries the segment, so if short, state the 16-times gap in one sentence and go to the next slide.}
\end{frame}

\begin{frame}{The roofline: bytes or arithmetic caps performance}
  \centering
  \mlsyssteps[width=0.92\textwidth,height=5.0cm,keepaspectratio]{roofline_step}{3}\\[0.1cm]
  \normalsize
  \only<1|handout:0>{Arithmetic intensity $=$ operations $\div$ bytes moved.}%
  \only<2|handout:0>{The chip cannot exceed its peak: 312 TFLOP/s.}%
  \only<3-|handout:1>{Ridge $= 312 \div 2.04 = 153$ FLOP/byte. Left of it, bytes limit; right of it, arithmetic.}
  \mlsyscredit{NVIDIA A100, FP16: System Paradigms, Napkin Math 2.1; roofline model after Williams, Waterman, and Patterson 2009.}
  \note{[3 min] This is the roofline, the first time it appears in the book; Lecture 11 develops it fully. Build it slowly. Explain the axes first: the horizontal axis is arithmetic intensity, the number of operations a workload performs for each byte it moves from memory; the vertical axis is the achieved rate, both on log scales. Step 1: if every byte arrives at 2.04 TB/s and each byte supports I operations, the rate is I times 2.04 TB/s. That is the blue slope. Step 2: the arithmetic units cannot exceed 312 TFLOP/s, the orange ceiling. Step 3: the two meet at the ridge point, peak divided by bandwidth, about 153 FLOP per byte for the A100. A workload left of the ridge is limited by memory bandwidth (the chapter says bandwidth-bound); right of it, by arithmetic (compute-bound). ENGAGE: ``A workload does 10 operations per byte. Which wall stops it?'' (Bandwidth; it reaches about 20 TFLOP/s, about 7 percent of peak.) WARN: ``memory-bound'' here means bandwidth, not running out of capacity; the book keeps the two words separate. FLEX: [CORE].}
\end{frame}

\begin{frame}{Single-image ResNet-50 sits just left of the A100's ridge}
  \begin{columns}[T]
    \begin{column}{0.47\textwidth}
      \begin{mlsyscard}{mlsysaccent}
        \textbf{Given}\\[0.05cm]
        ResNet-50, batch 1: 8.2 GFLOP\\
        Bytes moved (FP16): 56.3 MB\\
        A100: 312 TFLOP/s, 2.04 TB/s
      \end{mlsyscard}
      \footnotesize
      \uncover<2->{Compute:\enspace$8.2{\times}10^{9} \div 312{\times}10^{12} = 0.026$ ms\\[0.1cm]}
      \uncover<3->{Memory:\enspace$56.3{\times}10^{6} \div 2.04{\times}10^{12} = 0.028$ ms\\[0.1cm]}
      \uncover<4->{Intensity:\enspace$8.2{\times}10^{9} \div 56.3{\times}10^{6} = \mathbf{146}$, below 153}
    \end{column}
    \begin{column}{0.51\textwidth}
      \centering\uncover<4->{\safeimg[width=\textwidth,height=5.6cm,keepaspectratio]{roofline_resnet.pdf}}
    \end{column}
  \end{columns}
  \mlsyscredit{System Paradigms, Napkin Math 2.1. Peak rates, so both times are lower bounds.}
  \note{[3 min] A worked example on the roofline. ResNet-50 classifies one image with about 8.2 billion operations; with 16-bit weights and activations, the chapter's model moves about 56.3 MB. Step 2: compute time at peak, 0.026 ms. Step 3: memory time at full bandwidth, 0.028 ms. Memory is slightly slower, so the image is bandwidth-bound. Step 4: the same verdict from the roofline: 146 FLOP per byte, just left of the 153 ridge. With one image, each weight is read once and used once. On a flagship phone NPU (35 TOPS, 51.2 GB/s, INT8, 28.2 MB moved) compute takes 0.23 ms and memory 0.55 ms: bandwidth-bound by more. WARN: this is close to the ridge, so real kernel efficiencies decide which term wins in practice; the point is that batch 1 lives near or left of the ridge. ENGAGE: ``How would you move it right?'' (Process many images per weight read: batching.) FLEX: [CORE].}
\end{frame}

\begin{frame}{Predict: does 2$\times$ FLOP/s speed up text generation?}
  \mlsysactivity{Predict}{3 min}
  \begin{columns}[c]
    \begin{column}{0.56\textwidth}
      \centering
      \only<1|handout:0>{\safeimg[width=\textwidth,height=5.6cm,keepaspectratio]{roofline_decode_q.pdf}}%
      \only<2-|handout:1>{\safeimg[width=\textwidth,height=5.6cm,keepaspectratio]{roofline_decode_a.pdf}}
    \end{column}
    \begin{column}{0.42\textwidth}
      \small One user, one token at a time: about 1 FLOP per byte.\\[0.35cm]
      \large
      \choice{A}{About 2$\times$ faster}\\[0.15cm]
      \choice{B}{About 1.4$\times$ faster}\\[0.15cm]
      \alt<2->{\rchoice{C}{About the same}}{\choice{C}{About the same}}
    \end{column}
  \end{columns}
  \mlsyscredit{Decode intensity: Model Serving, \S13.9 (16-bit weights, about 2 operations per 2 bytes). Batching arrow assumes weight traffic dominates.}
  \note{[3 min] Generating text with a large language model produces one token at a time, and each token reads every weight of the model once. With 16-bit weights, that is about two operations for two bytes: about 1 FLOP per byte. ENGAGE: vote, discuss with a neighbor for 30 seconds, vote again. Step 2: answer C. The red point sits far down the memory slope at about 2 TFLOP/s, under 1 percent of peak. The dotted 2x ceiling is far above it and changes nothing; only more bandwidth or fewer bytes (smaller or lower-precision weights) help. Then point to the gray arrow: serving many users at once reuses each weight read for every user in the batch, so intensity rises roughly with batch size and the point slides up the slope toward the ceiling. WARN: the arrow ignores the per-user memory (the KV cache) that grows with batch, so the real crossover comes later; Lecture 13 does that accounting. This is why GPT-2 and Llama decode are the book's bandwidth-bound lighthouse. FLEX: [CORE].}
\end{frame}

\begin{frame}{Diagnose: which resource binds each workload?}
  \mlsysactivity{Diagnose}{2 min}
  \centering\footnotesize
  \renewcommand{\arraystretch}{1.1}
  \begin{tabular}{@{}P{0.23\textwidth}P{0.18\textwidth}P{0.35\textwidth}P{0.13\textwidth}@{}}
    \toprule
    \textbf{Workload} & \textbf{Bound by} & \textbf{Why} & \textbf{Home tier} \\
    \midrule
    ResNet-50 training & \leavevmode\uncover<2->{\textcolor{OrangeLine}{\textbf{Compute}}} & 8.2 GFLOP per image, large batches & Cloud \\
    GPT-2 / Llama decode & \leavevmode\uncover<2->{\textcolor{BlueLine}{\textbf{Bandwidth}}} & 1.5B to 70B weights read per token & Cloud \\
    DLRM recommendation & \leavevmode\uncover<2->{\textcolor{VioletLine}{\textbf{Capacity}}} & Embedding tables over 100 GB & Cloud \\
    MobileNetV2 & \leavevmode\uncover<2->{\textcolor{OrangeLine}{\textbf{Compute}}, efficient} & 13.7$\times$ fewer FLOPs than ResNet-50 & Mobile, edge \\
    Keyword spotting & \leavevmode\uncover<2->{\textcolor{RedLine}{\textbf{Power}}} & About 200 KB, always on & TinyML \\
    \midrule
    \leavevmode\uncover<3->{ResNet-50, batch 1} & \leavevmode\uncover<3->{\textcolor{BlueLine}{\textbf{Bandwidth}}} & \leavevmode\uncover<3->{146 FLOP/byte, below the 153 ridge} & \leavevmode\uncover<3->{Any} \\
    \bottomrule
  \end{tabular}\\[0.3cm]
  \uncover<3->{\normalsize The binding resource belongs to the \textbf{regime}, not the model.}
  \mlsyscredit{System Paradigms, \S2.3.2 and table 2; ResNet-50 at batch 1: Napkin Math 2.1.}
  \note{[3 min] The five lighthouse models from Lectures 0 and 1, now placed by what binds them first; the chapter calls these workload archetypes. ENGAGE: the Bound by column is hidden. Give pairs one minute to name the binding resource for each row from the Why column, choosing among compute, bandwidth, capacity, and power. Step 2: reveal. Two are terms of the iron law: compute-bound (ResNet-50 training) and bandwidth-bound (decode). Two are preconditions: capacity-bound (DLRM's embedding tables exceed one accelerator's memory) and power-bound (keyword spotting must run always on, on milliwatts). MobileNetV2 does the same task as ResNet-50 with 13.7 times fewer operations, sized for a phone's 3 to 5 W. A pair that answered compute for decode reasoned from model size; ask how many operations each byte supports (about one). Step 3: the same ResNet-50 served one image at a time sits at 146 FLOP per byte, below the 153 ridge, so it is bandwidth-bound, as the worked example showed. The archetype is a regime (model, batch size, and machine together), not a property of the model. LINK: Lecture 6 opens up each architecture. FLEX: [CORE].}
\end{frame}

% =============================================================================
% SECTION: POWER
% =============================================================================
\mlsysnote{[0.5 min] The third wall: power, and the heat it becomes. This is what separates a phone and a sensor from a server.}
\mlsysdivider{Why can a phone not run what a server runs, all day?}{System Paradigms, \S2.2, \S2.7, and \S2.8}
\section{Power}

\begin{frame}{Clock speeds stalled when power density hit its limit}
  \begin{columns}[c]
    \begin{column}{0.58\textwidth}
      \centering\safeimg[width=\textwidth,height=5.6cm,keepaspectratio]{microprocessor_trends.pdf}
    \end{column}
    \begin{column}{0.40\textwidth}
      \normalsize
      $P \propto C\,V^2 f$, and $V$ rises with $f$:\\[0.2cm]
      {\Large$P \propto f^3$}\\[0.4cm]
      \uncover<2->{{\Large\textbf{\textcolor{mlsysaccent}{2$\times$ clock $\approx$ 8$\times$ power}}}\\[0.4cm]}
      \uncover<3->{\small Transistors went into more cores and accelerators instead.}
    \end{column}
  \end{columns}
  \mlsyscredit{System Paradigms, equation 2 and fig.\ 2; data: Karl Rupp, Microprocessor Trend Data, 2022.}
  \note{[2 min] Real data first. Explain the plot: fifty years of processors on a log scale. Red is transistors, green clock frequency, blue single-thread performance, orange power. In the mid-2000s (the gray band) green flattens near 4.7 GHz and orange flattens too, while red keeps climbing about 190 times. Step 1: the equation. Dynamic power is capacitance times voltage squared times frequency; a faster clock needs a higher voltage, so power grows roughly with the cube of frequency. Step 2: doubling the clock costs about eight times the power, all of it heat. Step 3: until the mid-2000s, Dennard scaling held: shrinking transistors let voltage fall with size, so power density stayed constant and clocks could rise for free. When voltage stopped falling, that ended, and the extra transistors went into many cores, GPUs, and TPUs. This is why ML hardware is parallel and specialized. WARN: the cube assumes voltage scales with frequency; real chips vary, but the direction holds. FLEX: [OPTIONAL]; if short, skip it and state $P \propto f^3$ over the next slide.}
\end{frame}

\begin{frame}{A phone throttles once heat outruns its chassis}
  \centering
  \mlsyssteps[width=0.92\textwidth,height=4.6cm,keepaspectratio]{thermal_wall_step}{2}\\[0.1cm]
  \uncover<2->{\normalsize A thin, fanless phone can shed only about 3 W, whatever its battery holds.}
  \mlsyscredit{System Paradigms, Napkin Math 2.9 (the thermal wall). Illustrative trace.}
  \note{[2.5 min] The power wall at phone scale. If the previous slide was skipped, say first: dynamic power grows roughly with the cube of clock frequency, so heat, not transistors, caps speed. An unoptimized vision model draws 12 W; a fanless phone can remove only about 3 W of heat through its case. Step 1: the model starts at 100 frames per second and the chip warms at about 1 degree Celsius per second. Step 2: within about a minute it reaches its 80 degree trip point, the operating system lowers the clock, and throughput falls to about 30 frames per second, the rate 3 W can sustain. Benchmarks that run for ten seconds report the burst; users live on the sustained line. Suppose a redesign (lighter model, integer arithmetic) cuts power 4 times: 12 W becomes 3 W, exactly at the ceiling with no headroom. WARN: optimization can bring a model down to the ceiling; it cannot raise the ceiling. FLEX: [CORE].}
\end{frame}

\begin{frame}{Estimate: can an always-on detector last the day?}
  \mlsysactivity{Your turn}{3 min}
  \begin{columns}[T]
    \begin{column}{0.40\textwidth}
      \vspace{0pt}\begin{mlsyscard}{mlsysaccent}
        \textbf{Given}\\[0.1cm]
        Background object detector: 2 W\\
        Phone battery: 15 Wh\\
        Users expect: 24 hours
      \end{mlsyscard}
    \end{column}
    \begin{column}{0.52\textwidth}
      \vspace{0pt}\normalsize
      \uncover<2->{Runtime $= 15\ \text{Wh} \div 2\ \text{W} = 7.5$ hours\\[0.35cm]}
      \uncover<3->{A day needs $2\ \text{W} \times 24\ \text{h} = 48$ Wh\\[0.35cm]}
      \uncover<4->{{\Large\textbf{\textcolor{mlsysaccent}{320\% of the battery}}}\\[0.3cm]}
      \uncover<5->{Run it less often, or run a lighter model.}
    \end{column}
  \end{columns}
  \mlsyscredit{System Paradigms, Napkin Math 2.8 (the battery tax).}
  \note{[3 min] ENGAGE: one minute alone, then compare with a neighbor. Step 2: the feature alone empties the battery in 7.5 hours, before the screen, radio, or anything else runs. Step 3: a full day at 2 W needs 48 Wh. Step 4: 320 percent of the battery. Step 5: two fixes, duty cycling (run it only when needed) and a lighter model. Then contrast with the previous slide: the battery tax is a budget, so a bigger battery or less frequent use fixes it; the thermal wall is a ceiling on watts, so a bigger battery does nothing. Energy equals power times time. WARN: students often divide Wh by W and get confused by units; Wh divided by W is hours. FLEX: [CORE].}
\end{frame}

\begin{frame}{Energy per inference spans eight orders of magnitude}
  \centering
  \mlsyssteps[width=0.95\textwidth,height=5.0cm,keepaspectratio]{energy_per_inference_step}{2}\\[0.1cm]
  \uncover<2->{\normalsize Sending 1 MB costs 39 times the whole phone's inference energy.}
  \mlsyscredit{System Paradigms, fig.\ 7 and Napkin Math 2.5. Device-only energy; the LLM query is an order-of-magnitude anchor.}
  \note{[2.5 min] The book's figure; read the axis first: joules per inference, log scale, each grid line a factor of 100. A keyword spotter on dedicated low-power silicon uses about 10 microjoules; MobileNetV2 on a phone about 2.6 mJ; ResNet-50 on a Jetson about 15.9 mJ and on an A100 about 75.5 mJ; a cloud language-model query about 1 kJ. The span is about 100 million times. Step 2: the chapter's transmission example, the dashed red bar. Sending 1 MB over Wi-Fi or LTE costs about 100 mJ. The white tick inside the phone bar marks the NPU alone, about 0.1 mJ, so sending costs 1,000 times the NPU's work. The blue bar counts the whole phone during inference, 2.6 mJ, and sending still costs 39 times more. Say both numbers aloud so students do not read 0.1 and 2.6 as a contradiction: they draw different accounting boundaries. On a battery, sending data usually costs more energy than inferring on it. WARN: these bars use different tasks and accounting boundaries; they show scale, not a controlled comparison. FLEX: [CORE].}
\end{frame}

\begin{frame}{A tighter budget fails a model that fits in flash}
  \begin{columns}[c]
    \begin{column}{0.30\textwidth}
      \centering\safeimg[width=\textwidth,height=3.4cm,keepaspectratio]{esp32-s3-board.jpg}\\[0.1cm]
      {\footnotesize ESP32-S3, about \$10:\\512 KB SRAM, 4.2 MB flash}
    \end{column}
    \begin{column}{0.68\textwidth}
      \centering\mlsyssteps[width=\textwidth,height=4.6cm,keepaspectratio]{doorbell_step}{2}\\[0.1cm]
      \normalsize\only<1|handout:0>{200 ms chime: 101 ms per frame passes.}\only<2-|handout:1>{50 ms chime: the same model fails.}
    \end{column}
  \end{columns}
  \mlsyscredit{System Paradigms, \S2.3.2 and fig.\ 3 (smart doorbell); 51\% derived from 101 ms over 200 ms. Photo: VectorVoyager (CC BY-SA 4.0), Wikimedia Commons.}
  \note{[2 min] The bottom tier. A smart doorbell runs a visual wake-word model (is a person in view?) on this ESP32-S3. Explain the bars first: each is demand as a percent of its limit, and past the dashed budget line at 100 percent the design fails. Step 1: the 16-bit weights, 0.5 MB, fill only 11.9 percent of the 4.2 MB flash. The model needs 25 million operations per frame and the chip sustains about 0.25 GFLOP/s, so compute takes 100 ms; with about 1 ms of overhead, 101 ms per frame, 51 percent of a 200 ms chime. Step 2: a more responsive chime with a 50 ms budget. Nothing about the model changed, and latency now uses 202 percent of its budget. Flash was not the binding constraint; latency is. WARN: passing the flash check proves little; the activations must also fit in the 512 KB of SRAM, with no virtual memory to fall back on. That is why TinyML is not ``mobile, scaled down'': models are redesigned, not just quantized. FLEX: [OPTIONAL].}
\end{frame}

% =============================================================================
% SECTION: CHOOSING
% =============================================================================
\mlsysnote{[0.5 min] We have the walls. Now use them to decide where a model runs, and when one tier is not enough.}
\mlsysdivider{How do we pick a tier, and when do we use several?}{System Paradigms, \S2.9, \S2.10, and \S2.12}
\section{Choosing}

\begin{frame}{Screen for feasibility before comparing cost}
  \begin{columns}[c]
    \begin{column}{0.57\textwidth}
      \centering
      \begin{tikzpicture}[gate/.style={draw=midgray!70, fill=vlightgray!25, rounded corners=2pt, minimum width=8.4cm, minimum height=0.62cm, align=left, text width=8.0cm, font=\small, text=darktext}, hot/.style={gate, draw=RedLine, line width=1.4pt, fill=RedL!60}]
        \node[gate, dashed, text=midgray] (b) at (0,0.95) {\textbf{0. Baseline}: does ML beat a simple rule?};
        \node[gate] (p) at (0,0) {\textbf{1. Privacy}: may the data leave?};
        \node[gate] (l) at (0,-0.95) {\textbf{2. Latency}: does a remote path fit the budget?};
        \node<2->[hot] at (0,-0.95) {\textbf{2. Latency}: does a remote path fit the budget?};
        \node[gate] (r) at (0,-1.9) {\textbf{3. Resources}: do memory, compute, energy fit?};
        \node[gate] (c) at (0,-2.85) {\textbf{4. Cost}: cheapest survivor over its lifetime};
        \foreach \a/\b in {b/p,p/l,l/r,r/c} \draw[-{Stealth}, midgray, line width=0.8pt] (\a.south) -- (\b.north);
        \draw[-{Stealth}, RedLine, line width=1.1pt] (c.south) -- ++(0,-0.45) node[below, font=\small, text=RedLine, align=center] {\textbf{No survivor}: redesign, or split across tiers};
      \end{tikzpicture}
    \end{column}
    \begin{column}{0.40\textwidth}
      \normalsize Emergency braking at 100 km/h:\\[0.2cm]
      \small 100 ms budget $=$ 2.8 m of travel.\\
      Cloud round trip: 50 to 150 ms.\\[0.3cm]
      \uncover<2->{\normalsize The latency gate removes the cloud before cost is asked.}
    \end{column}
  \end{columns}
  \mlsyscredit{After System Paradigms, fig.\ 11; Example 2.1.}
  \note{[2.5 min] Gate 0, in gray, comes before any tier: does ML beat a simple rule at all? The chapter calls the extra operating cost of an ML system (data pipelines, monitoring, retraining) the complexity tax; Lecture 3 starts there. Then the four screening gates, in order; each passes only its survivors to the next. Privacy goes first because a legal rule removes an option without any measurement. Walk one example. A car's pedestrian detector must decide within 100 ms; at 100 km/h the car covers 2.8 m in that time. A cloud round trip of 50 to 150 ms uses half the budget at best and all of it at worst. Step 2: the latency gate removes the cloud; an automotive edge accelerator remains and must still pass the resource gate. The red exit: if no tier survives, the answer is not the cheapest failure; change the workload, or split it across tiers, which the next three slides do. WARN: cost is compared last and only among feasible options, over the whole lifetime (total cost of ownership). In the chapter's cloud-against-edge comparison (Napkin Math 2.3), labor, not hardware, is about 69 percent of the edge budget. FLEX: [CORE].}
\end{frame}

\begin{frame}{Discuss: where should each voice stage run?}
  \mlsysactivity{Discuss}{2.5 min}
  \begin{columns}[c]
    \begin{column}{0.30\textwidth}
      \centering\mlsysphoto[width=\textwidth,height=4.0cm,keepaspectratio]{amazon-echo-dot-with-mug.jpg}
    \end{column}
    \begin{column}{0.68\textwidth}
      \centering
      \begin{tikzpicture}[st/.style={draw=midgray!70, fill=vlightgray!25, rounded corners=2pt, minimum width=2.75cm, minimum height=1.1cm, align=center, font=\small, text=darktext}]
        \node[st] (w) at (0,0) {Wake word};
        \node[st] (s) at (3.4,0) {Speech\\to text};
        \node[st] (u) at (6.8,0) {Language\\understanding};
        \draw[-{Stealth}, midgray, line width=0.9pt] (w) -- (s);
        \draw[-{Stealth}, midgray, line width=0.9pt] (s) -- (u);
        \node<1|handout:0>[font=\normalsize\bfseries] at (3.4,-1.05) {Which tier runs each?};
        \node<2->[font=\small, text=RedLine, align=center] at (0,-1.05) {\textbf{TinyML}\\always on};
        \node<2->[font=\small, text=OrangeLine, align=center] at (3.4,-1.05) {\textbf{Mobile}\\private, offline};
        \node<2->[font=\small, text=BlueLine, align=center] at (6.8,-1.05) {\textbf{Cloud}\\large model};
      \end{tikzpicture}\\[0.5cm]
      {\normalsize\textcolor{midgray}{Two minutes with a neighbor.}}
    \end{column}
  \end{columns}
  \mlsyscredit{Photo: Pittigrilli (CC BY-SA 4.0), Wikimedia Commons. Stages: System Paradigms, \S2.10.}
  \note{[2.5 min] Lecture 0 asked where the wake word should run; now place the whole pipeline. ENGAGE: pairs assign a tier to each of the three stages and give one reason each. Take three answers. Step 2: the chapter's answer. The wake word runs on a microcontroller-class core because it listens all day on milliwatts. Speech to text can run on the phone, which keeps raw audio private and works offline. Language understanding needs a model larger than the phone holds, so it runs in the cloud. Accept other defensible answers, such as cloud speech recognition on a low-end speaker; what matters is the reason. The next slide shows why the first stage cannot be in the cloud at scale. FLEX: [CORE].}
\end{frame}

\begin{frame}{Detecting the wake word on the device removes 32 TB/s}
  \centering
  {\small\textcolor{midgray}{1 billion devices $\cdot$ 16 kHz, 16-bit audio $=$ 32 KB/s each $\cdot$ 20 queries a day each}}\\[0.15cm]
  \mlsyssteps[width=0.9\textwidth,height=3.0cm,keepaspectratio]{voice_upload_step}{2}\\[0.25cm]
  \normalsize
  \uncover<3->{Streaming everything to the cloud: about \$500M a year.\\[0.15cm]}
  \uncover<4->{After the split: 1.1M GPU-hours a day, about 20.8 datacenters at peak.}
    \mlsyscredit{System Paradigms, Napkin Math 2.10 (the voice assistant wall). 37 GB/s derived from 32 TB/s and the chapter's 860$\times$.}
  \note{[2 min] The chapter's capacity scenario. Step 1: if wake-word detection ran in the cloud, every microphone would stream 16 kHz, 16-bit audio, 32 KB/s, so one billion devices send 32 TB/s, all day. Step 2: with the wake word detected on the device, each device uploads only about five seconds per query, cutting audio leaving the device about 860 times, to roughly 37 GB/s in total. Step 3: the chapter puts cloud-side detection at about \$0.50 per device per year, \$500 million a year for a free feature; the local detector costs under a cent per device per year at 0.1 to 1 mW. Step 4, [OPTIONAL]; skip it if short: the cloud still has real work. 20 billion queries a day at about 200 ms of GPU time each is 1.1 million GPU-hours a day; with 10,000 GPUs per datacenter that is 4.6 datacenters on average and, with a 4.5 times peak, about 20.8 at peak. The split relieves the battery and the datacenter at once. FLEX: [CORE].}
\end{frame}

\begin{frame}[t]{Three patterns split one pipeline across tiers}
  \centering
  \vspace*{-0.25cm}
  \begin{tikzpicture}[
      tier/.style={draw=midgray!45, fill=vlightgray!30, line width=0.6pt, rounded corners=2pt},
      tt/.style={font=\footnotesize\bfseries, text=darktext},
      fl/.style={line width=1pt, midgray!60, -{Stealth[length=5pt]}},
      hl/.style={line width=2pt, mlsysaccent, -{Stealth[length=6pt]}},
      hb/.style={draw=mlsysaccent, line width=2pt},
      ann/.style={font=\footnotesize\bfseries, text=mlsysaccent, align=center},
      tag/.style={fill=white, inner sep=1pt, font=\footnotesize, text=darktext}]
    % tier containers (drawn first, behind the stages)
    \draw[tier] (-5.45,1.47) rectangle (3.35,2.20);
    \draw[tier] (-6.55,1.10) rectangle (-3.85,-1.16);
    \draw[tier] (3.85,1.10) rectangle (6.55,-1.16);
    \draw[tier] (-1.35,-0.05) rectangle (1.35,-2.28);
    \node[tt] at (-4.45,1.84) {Cloud ML};
    \node[tt] at (-5.20,-0.94) {Edge ML};
    \node[tt] at (5.20,-0.94) {Mobile ML};
    \node[tt] at (0,-2.06) {TinyML};
    % stages
    \node[bx=Red] (ct) at (-2.0,1.84) {Training};
    \node[bx=Red] (ca) at (2.0,1.84) {Analytics};
    \node[bx=Blue] (ei) at (-5.2,0.75) {Inference};
    \node[bx=Blue] (ep) at (-5.2,-0.40) {Processing};
    \node[bx=Orange] (mi) at (5.2,0.75) {Inference};
    \node[bx=Orange] (mp) at (5.2,-0.40) {Processing};
    \node[bx=Green] (ti) at (0,-0.40) {Inference};
    \node[bx=Green] (ts) at (0,-1.52) {Sensors};
    % within each tier
    \draw[fl] (ei) -- node[tag, inner ysep=0.5pt] {Results} (ep);
    \draw[fl] (mi) -- node[tag, inner ysep=0.5pt] {Results} (mp);
    \draw[fl] (ts) -- node[tag, inner ysep=0.5pt] {Data} (ti);
    % deploy: cloud training down to every inference stage
    \draw[fl] (ct.south) -- (-2.0,1.25) -| (-4.7,1.01);
    \draw[fl] (-2.0,1.25) -| (5.7,1.01);
    \draw[fl] (-2.0,1.25) -- (0,1.25) -- (ti.north);
    % assist: mobile inference hands off to edge inference
    \draw[fl] (4.7,1.01) -- (4.7,2.78) -- (-5.7,2.78) -- (-5.7,1.01);
    % results up from edge processing to cloud analytics
    \draw[fl] (ep.west) -- (-6.95,-0.40) -- (-6.95,2.42) -- (2.0,2.42) -- (ca.north);
    % data up from mobile processing to cloud analytics
    \draw[fl] (mp.east) -- (6.95,-0.40) -- (6.95,1.84) -- (ca.east);
    % TinyML to its neighbors
    \draw[fl] (ti.west) -- (ep.east);
    \draw[fl] (ti.east) -- (mp.west);
    % pattern 1, train-serve split: the three Deploy arrows
    \begin{scope}[visible on={<1,4|handout:1>}]
      \draw[hl] (ct.south) -- (-2.0,1.25) -| (-4.7,1.01);
      \draw[hl] (-2.0,1.25) -| (5.7,1.01);
      \draw[hl] (-2.0,1.25) -- (0,1.25) -- (ti.north);
      \node[ann] at (-2.3,0.62) {1: train 16 B/param,\\serve 2 B/param};
    \end{scope}
    % pattern 2, hierarchical processing: results and data flow upward
    \begin{scope}[visible on={<2,4|handout:1>}]
      \draw[hl] (ep.west) -- (-6.95,-0.40) -- (-6.95,2.42) -- (2.0,2.42) -- (ca.north);
      \draw[hl] (mp.east) -- (6.95,-0.40) -- (6.95,1.84) -- (ca.east);
      \draw[hl] (ti.west) -- (ep.east);
      \node[ann] at (2.1,0.62) {2: voice assistant,\\factory cameras};
    \end{scope}
    % pattern 3, progressive deployment: one model family, one size per tier
    \begin{scope}[visible on={<3,4|handout:1>}]
      \node[bx=Blue, hb] at (ei) {Inference};
      \node[bx=Orange, hb] at (mi) {Inference};
      \node[bx=Green, hb] at (ti) {Inference};
      \node[ann] at (-4.4,-1.72) {3: one model family,\\one size per tier};
    \end{scope}
    % flow labels last, so highlights never cover them
    \node[tag] at (-3.05,1.25) {Deploy};
    \node[tag] at (0,2.78) {Assist};
    \node[tag] at (-2.9,2.42) {Results};
    \node[tag] at (6.95,0.72) {Data};
    \node[tag] at (-2.6,-0.40) {Results};
    \node[tag] at (2.6,-0.40) {Sync};
  \end{tikzpicture}\\[0.1cm]
  \small\textbf{\alert<1|handout:0>{1\enspace Train-serve split}\quad$\cdot$\quad\alert<2|handout:0>{2\enspace Hierarchical processing}\quad$\cdot$\quad\alert<3|handout:0>{3\enspace Progressive deployment}}
  \mlsyscredit{After System Paradigms, fig.\ 12; patterns from table 10; 16 B/param: Rajbhandari et al.\ 2020, as cited in \S2.10.}
  \note{[3 min] The book's figure of how tiers talk. Orient first: four tier boxes; models deploy downward, results and data flow upward. Two arrows are not part of the three patterns. Assist: a phone hands a request it cannot serve to a nearby edge box. Sync: devices exchange updates with each other. Step 1, train-serve split: the three Deploy arrows. Training keeps weights, gradients, and optimizer state, about 16 bytes per parameter, against 2 bytes for 16-bit inference, about 8 times more, so training stays in the cloud and only the trained weights travel down. Step 2, hierarchical processing: the upward Results and Data arrows. Each tier sends up only what it found; the voice assistant (wake word, then phone, then cloud) and the factory cameras (defect records, not video) are both instances. Step 3, progressive deployment: the three Inference boxes run one model family, one size per tier, sharing training data and labels. Step 4 shows all three together. ENGAGE: ``Which pattern is the robot arm's design?'' (Train-serve split: the cloud trains, the edge box beside the arm infers.) WARN: hybrid is not ``running two models''; the tiers must share feature definitions, preprocessing, and model versions, or they silently disagree. LINK: Lecture 14. FLEX: [CORE].}
\end{frame}

\begin{frame}{Check: true or false?}
  \mlsysactivity{Vote}{2 min}
  \normalsize
  \begin{tabular}{@{}P{0.70\textwidth}c@{}}
    \toprule
    \textbf{Claim} & \textbf{Verdict} \\
    \midrule
    A 10$\times$ faster model makes a 200 ms camera pipeline 10$\times$ faster. & \leavevmode\uncover<2->{\textbf{\textcolor{RedLine}{False: 1.37$\times$}}} \\ \addlinespace[0.3cm]
    Quantizing to INT8 lets a 12 W model run all day on a 3 W phone. & \uncover<3->{\textbf{\textcolor{RedLine}{False}}} \\ \addlinespace[0.3cm]
    Faster 5G links will let a 10 ms control loop use a cloud 1,500 km away. & \uncover<4->{\textbf{\textcolor{RedLine}{False: 15 ms}}} \\
    \bottomrule
  \end{tabular}
  \mlsyscredit{System Paradigms, \S2.12 (fallacies and pitfalls).}
  \note{[2 min] The chapter's fallacies as votes; thumbs up or down for each before any reveal. Step 2: the phone camera path is 100 ms of image processing, 60 ms of ML, and 40 ms of tone mapping. A 10x faster model makes 60 ms into 6 ms, so 200 ms becomes 146 ms, 1.37 times faster; even an infinitely fast model gives at most 1.43 times (Amdahl's law; Lecture 11). Step 3: quantization shrinks weights about 4 times but does not cut power in proportion, because switching and leakage remain; the 3 W ceiling holds. Step 4: 5G raises bandwidth, not the speed of light. Propagation alone over 1,500 km and back is 15 ms, past the 10 ms loop before any compute; this is the arm from the opening, and the chapter's self-check asks the same question. FLEX: [CORE].}
\end{frame}

% =============================================================================
% SYNTHESIS AND CLOSE
% =============================================================================
\begin{frame}{Back to the arm: the regime decided the design}
  \begin{columns}[c]
    \begin{column}{0.44\textwidth}
      \centering\safeimg[width=\textwidth,height=4.6cm,keepaspectratio]{robot-arms-bmw-leipzig.jpg}
    \end{column}
    \begin{column}{0.53\textwidth}
      \normalsize
      \uncover<1->{\textcolor{RedLine}{\textbf{Cloud}}: 15 ms round trip $>$ 10 ms. Rejected.\\[0.35cm]}
      \uncover<2->{\textcolor{GreenLine}{\textbf{Edge box}} beside the arm: no propagation. Measure the local path.\\[0.35cm]}
      \uncover<3->{\textcolor{BlueLine}{\textbf{Cloud}} still trains the model: a train-serve split.}
    \end{column}
  \end{columns}
  \mlsyscredit{Photo: BMW Werk Leipzig (CC BY-SA 2.0 de), Wikimedia Commons. System Paradigms, \S2.5, \S2.6, \S2.10.}
  \note{[2 min] Return to the opening question with everything from today. Step 1: the light barrier rejects the cloud path; the latency gate. Step 2: an edge computer beside the arm removes propagation; it must still pass the resource gate (the detector must fit and sustain its frame rate within its power budget) and the complete path must be measured. Step 3: the cloud keeps the job it is good at, training and retraining on data from many lines: a hybrid design. ENGAGE: ``What would change if the arm needed only a 1-second response?'' (The cloud passes the latency gate, and cost decides.) Distance, not model speed, decided where the arm's model runs. FLEX: [CORE].}
\end{frame}

\begin{frame}{Takeaways}
  \small
  \renewcommand{\arraystretch}{1.0}
  \begin{tabular}{@{}>{\centering\arraybackslash}m{3.1cm}@{\hspace{0.4cm}}>{\raggedright\arraybackslash}m{0.72\textwidth}@{}}
    \thumb{power_tiers_step4.pdf} & Four regimes span nine orders of magnitude in power; each tier binds a different constraint. \\[0.12cm]
    \thumb{rtt_distance_step2.pdf} & Light covers 200 km of fiber per millisecond; when sending the data misses the deadline, the computation moves to the data. \\[0.12cm]
    \thumb{roofline_resnet.pdf} & Below the ridge point, bytes limit speed; batch-1 inference sits on the memory slope. \\[0.12cm]
    \thumb{thermal_wall_step2.pdf} & On battery-powered devices, energy per inference and a few watts of heat decide what runs. \\[0.12cm]
    \thumb{voice_upload_step2.pdf} & Screen feasibility first, compare cost last, and split the pipeline across tiers when no single tier survives. \\
  \end{tabular}
  \note{[1.5 min] One line per segment, each beside the figure that taught it; ask students to circle the one they could not yet explain to a friend and reread that section. One more idea from the chapter's close: the same techniques serve every tier, because the walls are the same; moving fewer bytes helps a microcontroller and a datacenter alike. That is why the rest of the course applies everywhere.}
\end{frame}

\begin{frame}{Open questions}
  \normalsize
  \begin{itemize}\setlength{\itemsep}{0.4cm}
    \item As phones gain memory and faster NPUs, which workloads will leave the cloud, and which never can?
    \item Who should decide where a model runs: the ML engineer, the hardware team, or the privacy lawyer?
    \item If round trips cannot shrink below physics, what hybrid designs will large models need?
  \end{itemize}
  \note{[1 min] No single answer. Assign one as a short written reflection, or use one to open the next discussion section.}
\end{frame}

\begin{frame}{Read and explore}
  \small
  \begin{itemize}\setlength{\itemsep}{0.2cm}
    \item \textbf{Read:} System Paradigms, \S2.1 to \S2.10 and \S2.12; redo Napkin Math 2.1, 2.4, 2.8, and 2.10 by hand.
    \item Williams, Waterman, and Patterson, ``Roofline: An Insightful Visual Performance Model,'' CACM 2009.
    \item Wulf and McKee, ``Hitting the Memory Wall: Implications of the Obvious,'' 1995.
    \item Hennessy and Patterson, ``A New Golden Age for Computer Architecture,'' CACM 2019.
    \item Kang et al., ``Neurosurgeon: Collaborative Intelligence Between the Cloud and Mobile Edge,'' ASPLOS 2017.
  \end{itemize}
  \note{[1 min] The chapter first, with its self-check questions. Each outside reading backs one segment: the roofline paper (Memory), the memory wall (Memory), the end of Dennard scaling and the turn to accelerators (Power), and splitting one network between phone and cloud (Choosing).}
\end{frame}

\begin{frame}[c]
  \centering
  {\small\textcolor{lightgray}{NEXT LECTURE \textbar{} THE ML LIFECYCLE}}\\[0.5cm]
  {\LARGE\bfseries Why is seeing the whole map necessary\\before walking any single path?}
  \note{[0.5 min] Today said where a model can run. Lecture 3 asks when a team must know it: a deployment budget discovered after five months of work invalidates those months. Read the ML Lifecycle chapter before then.}
\end{frame}

\end{document}
